\begin{figure}[ht]
\centering
\begin{tikzpicture}[x=1.0cm, y=0.9cm, font=\small]

% Axes: horizontal = risk (std dev), vertical = return.
\draw[->, thick] (-0.2, 0) -- (8.6, 0)
  node[right, font=\scriptsize] {risk $\sqrt{\vect{w}^{\transpose}\Sigma\vect{w}}$};
\draw[->, thick] (0, -0.2) -- (0, 6.4)
  node[above, font=\scriptsize] {return $\boldsymbol{\mu}^{\transpose}\vect{w}$};

% Cloud of random feasible portfolios (interior of the achievable region).
\foreach \x/\y in {%
  3.1/1.4, 3.6/1.1, 4.2/2.0, 4.8/1.6, 5.3/2.4, 5.9/2.0,
  3.4/2.2, 4.0/1.3, 4.6/2.7, 5.1/1.9, 5.7/3.0, 6.3/2.6,
  4.3/3.1, 4.9/2.3, 5.5/3.5, 6.0/3.1, 6.6/3.6, 7.0/3.2,
  3.8/2.6, 5.0/3.4, 5.8/2.7, 6.4/4.0, 7.2/4.1, 6.8/3.0} {
  \fill[enggray!55] (\x, \y) circle (1.0pt);
}

% The efficient frontier: upper-left convex boundary (parabola opening right).
% Parametrise risk r in [2.6, 7.6], return = 1.4 + 1.55*sqrt(r-2.6).
\draw[engblue, very thick, smooth, samples=40, domain=2.6:7.6]
  plot ({\x}, {1.4 + 1.55*sqrt(\x - 2.6)});
\node[engblue, anchor=west, font=\scriptsize] at (5.6, 5.4)
  {efficient frontier};

% The lower (inefficient) boundary for completeness.
\draw[enggray, thin, dashed, smooth, samples=40, domain=2.6:7.6]
  plot ({\x}, {1.4 - 0.55*sqrt(\x - 2.6)});

% Minimum-variance portfolio (leftmost point of the region).
\fill[engred] (2.6, 1.4) circle (2.2pt);
\node[engred, anchor=east, font=\scriptsize] at (2.45, 1.4)
  {min-variance};

% Three marked portfolios along the frontier for gamma = 10, 1, 0.1.
\fill[engorange] (3.6, 2.95) circle (2.0pt);
\node[engorange, anchor=north west, font=\scriptsize] at (3.65, 2.85)
  {$\gamma=10$};
\fill[engorange] (5.0, 3.92) circle (2.0pt);
\node[engorange, anchor=north west, font=\scriptsize] at (5.05, 3.82)
  {$\gamma=1$};
\fill[engorange] (7.0, 4.86) circle (2.0pt);
\node[engorange, anchor=north east, font=\scriptsize] at (6.9, 4.96)
  {$\gamma=0.1$};

\end{tikzpicture}
\caption[The Markowitz efficient frontier]{The achievable
risk-return region for a portfolio of risky assets. Each grey dot is
a feasible portfolio $\vect{w}$ on the simplex; the cloud fills a
region whose upper-left boundary is the efficient frontier (solid).
A portfolio below the frontier is dominated: another portfolio
offers the same return at lower risk, or more return at the same
risk. The minimum-variance portfolio (red) is the leftmost point.
Sweeping the risk-aversion parameter $\gamma$ from large to small
traces the frontier upward and to the right: a risk-averse investor
($\gamma=10$) sits near the minimum-variance point; a
risk-tolerant investor ($\gamma=0.1$) concentrates in
high-return, high-variance assets. Each marked point is the
solution of one convex quadratic program.}
\label{fig:vol02:ch15:efficient-frontier}
\end{figure}
